\author{}
\documentclass[11pt,a4paper]{article}
\usepackage[utf8]{inputenc}
\usepackage[T1]{fontenc}
\usepackage[french]{babel}

% Paquets pour les mathématiques
\usepackage{amsmath, amssymb, array, physics}

% Marges
\usepackage{geometry}
\geometry{hmargin=2cm, vmargin=2cm}

% Paquet pour les boîtes colorées
\usepackage[many]{tcolorbox}

% Définition d'un style personnalisé pour les définitions (bleu)
\newtcolorbox{boiteDef}[1]{
    colback=blue!5!white,
    colframe=blue!75!black,
    title={\textbf{#1}},
    arc=3mm,
    boxrule=1pt,
    fonttitle=\bfseries
}

% Définition d'un style pour les propriétés (rose)
\newtcolorbox{boiteProp}[1]{
    colback=pink!5!white,
    colframe=purple!70!pink,
    title={\textbf{#1}},
    arc=3mm,
    boxrule=1pt,
}

% Définition d'un style pour les rappels (vert)
\newtcolorbox{boiteRappel}[1]{
    colback=green!10!white,
    colframe=green!50!black,
    title={\textbf{#1}},
    arc=3mm,
    boxrule=1pt,
}

% Définition d'un style pour les remarques (jaune)
\newtcolorbox{boiteRemarque}[1]{
    colback=yellow!10!white,
    colframe=yellow!70!black,
    title={\textbf{#1}},
    arc=3mm,
    boxrule=1pt,
}

\renewcommand{\thesection}{\Roman{section}}

\begin{document}

\title{\textbf{Chapitre 12 : Primitives, équations différentielles}}
\date{}
\maketitle

\section{Equation différentielle $y' = f$ et primitives}

\subsection{Lien entre équation différentielle y' = f et primitives}

\begin{boiteRemarque}{Vocabulaire :}
    \begin{itemize}
        \item Une équation différentielle est une équation où l'inconnue est une fonction et où interviennent des dérivées de cette fonction.
        \item Résoudre une équation différentielle, c'est trouver toutes les fonctions solutions de l'équation.
    \end{itemize}
\end{boiteRemarque}

\begin{boiteDef}{Définition :}
    Soit $f$ une fonction définie sur un intervalle $I$. 
    \newline
    On dit que $F$ est une primitive de la fonction $f$ sur $I$ si $F$ est dérivable sur l'intervalle $I$ et que pour tout $x$ de $I$, $F'(x) = f(x)$.
\end{boiteDef}

\begin{boiteRemarque}{Remarques :}
    \begin{itemize}
        \item Une primitive de la fonction $f$ sur l'intervalle $I$ est une solution de l'équation différentielle $y' = f$.
        \item Pour une même fonction, il y a une infinité de primitives.
        \item $F$ est \textbf{une} primitive de la fonction $f$ sur $I$ signifie que la fonction $f$ est \textbf{la} fonction dérivée de $F$ sur $I$.
        \item La fonction $f'$ est \textbf{la} fonction dérivée de $f$ sur $I$ signifie que la fonction $f$ est \textbf{une} primitive de $f'$ sur $I$.
    \end{itemize}   
\end{boiteRemarque}

\begin{boiteProp}{Propriété :}
    Toute fonction $f$ continue sur un intervalle $I$ admet des primitives sur $I$.
\end{boiteProp}

\subsection{Primitives d'une même fonction}

\begin{boiteProp}{Propriétés :}
    Soit $F$ une primitive de la fonction $f$ sur l'intervalle $I$.
    \begin{itemize}
        \item Pour tout réel $k$, la fonction $G : x \mapsto F(x) + k$ est une primitive de la fonction $f$ sur $I$.
        \item Toute primitive de $f$ sur $I$ est de ce type.
    \end{itemize}
\end{boiteProp}

\begin{boiteRemarque}{Remarque :}
    Si la fonction $H$ est aussi une primitive de la fonction $f$ sur l'intervalle $I$, alors il existe un réel $k$ tel que pour tout $x$ de $I$, $H(x) - F(x) = k$.
\end{boiteRemarque}

\begin{boiteRemarque}{Vocabulaire}
    On dit que deux primitives d'une même fonction continue sur un intervalle diffèrent d'une constante
    \newline
    L'équation différentielle $y' = f$ a pour solutions toutes les primitives de $f$ sur $I$.    
\end{boiteRemarque}

\begin{boiteProp}{Propriété :}
    Soit $f$ une fonction continue sur $I$.
    \newline
    Il existe une unique primitive $G$ de la fonction $f$ sur l'intervalle $I$ telle que : $G(x_0) = y_0$
    \newline
    où $x_0$ est un nombre réel donné de $I$ et $y_0$ un nombre réel donné.
\end{boiteProp}

\section{Formulaire de primitives}

\subsection{Primitives de $f + g$ et de $\lambda f$}

\begin{boiteProp}{Propriété :}
    Soit $f$ et $g$ deux fonctions continues sur un intervalle $I$ de primitives respectives $F$ et $G$ sur l'intervalle $I$.
    \begin{itemize}
        \item La fonction $(F + G)$ est une primitive de la fonction $(f + g)$ sur l'intervalle $I$.
        \item $\forall \lambda \in \mathbb{R}$, la fonction $(\lambda F)$ est une primitive de la fonction $(\lambda f)$ sur l'intervalle $I$.
    \end{itemize}
\end{boiteProp}

\begin{boiteRemarque}{Remarque :}
    Aucune propriété sur $F \times G$ ou sur $\frac{F}{G}$.
\end{boiteRemarque}

\subsection{Primitives des fonctions de référence}

On lit le tableau des dérivées dans l'autres sens.
\newline

\newcolumntype{C}[1]{>{\centering\arraybackslash}p{#1}}
\renewcommand{\arraystretch}{0.7} % 1.2, 1.5, etc.
\begin{tabular}{|C{4cm}|C{7cm}|C{4cm}|}
\hline
\textbf{Fonction} & \textbf{Primitives} & \textbf{Intervalle $I$} \\
\hline
\color{black}{$x \mapsto a$} 
& \color{black}{$x \mapsto ax + b,\; b \in \mathbb{R}$} 
& \color{black}{$\mathbb{R}$} \\
\hline
\color{black}{$x \mapsto x$} 
& \color{black}{$x \mapsto \dfrac{1}{2} x^2 + k,\; k \in \mathbb{R}$} 
& \color{black}{$\mathbb{R}$} \\
\hline
\color{black}{$x \mapsto x^2$} 
& \color{black}{$x \mapsto \dfrac{1}{3} x^3 + k,\; k \in \mathbb{R}$} 
& \color{black}{$\mathbb{R}$} \\
\hline
\color{black}{$x \mapsto x^n,\; n \in \mathbb{Z}^*,\; \newline n \neq -1$}
& \color{black}{$x \mapsto \dfrac{1}{n+1}x^{n+1} + k,\; k \in \mathbb{R}$}
& \color{black}{$\mathbb{R}$ si $n>0$} \\
\hline
\color{black}{$x \mapsto \dfrac{1}{x^2}$}
& \color{black}{$x \mapsto -\dfrac{1}{x} + k,\; k \in \mathbb{R}$}
& \color{black}{$]-\infty;0[ \text{ et } ]0;+\infty[$} \\
\hline
\color{black}{$x \mapsto \dfrac{1}{\sqrt{x}}$}
& \color{black}{$x \mapsto 2\sqrt{x} + k,\; k \in \mathbb{R}$}
& \color{black}{$]0;+\infty[$} \\
\hline
\color{black}{$x \mapsto \dfrac{1}{x}$}
& \color{black}{$x \mapsto \ln(x) + k,\; k \in \mathbb{R}$}
& \color{black}{$]0;+\infty[$} \\
\hline
\color{black}{$x \mapsto \dfrac{1}{x}$}
& \color{black}{$x \mapsto \ln(-x) + k,\; k \in \mathbb{R}$}
& \color{black}{$]-\infty;0[$} \\
\hline
\color{black}{$x \mapsto e^x$}
& \color{black}{$x \mapsto e^x + k,\; k \in \mathbb{R}$}
& \color{black}{$\mathbb{R}$} \\
\hline
\color{black}{$x \mapsto \cos(x)$}
& \color{black}{$x \mapsto \sin(x) + k,\; k \in \mathbb{R}$}
& \color{black}{$\mathbb{R}$} \\
\hline
\color{black}{$x \mapsto \sin(x)$}
& \color{black}{$x \mapsto -\cos(x) + k,\; k \in \mathbb{R}$}
& \color{black}{$\mathbb{R}$} \\
\hline
\end{tabular}

\subsection{Primitives et composition}

\begin{boiteProp}{Propriété :}
    Si $v$ est une fonction dérivable sur un intervalle $J$ et $u$ est une fonction dérivable sur un intervalle $I$ tel que pour tout réel $x \in I$, $u(x) \in J$, alors la fonction $u' \times (v' \circ u)$ admet pour primitive sur $I$ la fonction $v \circ u$.
\end{boiteProp}

On en déduit le tableau suivant :
\newline

\begin{tabular}{|C{4cm}|C{7cm}|C{4cm}|}
\hline
\textbf{Fonction $f$} 
& \textbf{Une primitive $F$} 
& \textbf{Conditions} \\
\hline
\color{black}{$u' u^{n}$}
& \color{black}{$\dfrac{1}{n+1} u^{n+1}$}
& \color{black}{$n \in \mathbb{Z}^*,\; n \neq -1$} \\
\color{black}{}
& \color{black}{}
& \color{black}{Pour $n < -1$,\, u(x) $\neq 0$ pour tout $x \in I$} \\
\hline
\color{black}{$\dfrac{u'}{u^{2}}$}
& \color{black}{$-\dfrac{1}{u}$}
& \color{black}{$u(x) \neq 0$ pour tout $x \in I$} \\
\hline
\color{black}{$\dfrac{u'}{\sqrt{u}}$}
& \color{black}{$2 \sqrt{u}$}
& \color{black}{Pour tout $x \in I$,\, $u(x) > 0$} \\
\hline
\color{black}{$u' e^{u}$}
& \color{black}{$e^{u}$}
& \color{black}{(aucune condition particulière)} \\
\hline
\color{black}{$\dfrac{u'}{u}$}
& \color{black}{$\ln(u)$}
& \color{black}{Pour tout $x \in I$,\, $u(x) > 0$} \\
\hline
\end{tabular}


\section{Equations différentielles}

\subsection{Equations différentielles du type $y' = ay$ où $a \in \mathbb{R^*}$}

\begin{boiteProp}{Propriété :}
    Les fonctions solutions de l'équation différentielle $y' = ay$ où $a \in \mathbb{R^*}$ sont les fonctions $f_k$ définies sur $\mathbb{R}$ par : $f_k(x) = k e^{ax}$ où $k \in \mathbb{R}$.
\end{boiteProp}

\subsection{Equation différentielle $y' = ay + b$}

\begin{boiteProp}{Propriété :}
    Les fonctions solutions de l'équation différentielle $y' = ay + b$ où $a$ et $b$ sont des réels non nuls sont les fonctions $f_k$ définies sur $\mathbb{R}$ par : $f_k(x) = k e^{ax} - \frac{b}{a}$ où $k \in \mathbb{R}$.
\end{boiteProp}

\subsection{Equation diférentielle $y' = ay + f$ ($f$ fonction, $a$ réel)}

\begin{boiteProp}{Propriété :}
    Les solutions de l'équation différentielle $y' = ay + f$ où $a$ est un réel non nul et $f$ une fonction sont les fonctions : $x \mapsto k e^{ax} + g(x)$
    \newline
    où $g$ est une solution de l'équation différentielle et $k$ une constante réelle.
\end{boiteProp}


















\end{document}
